% =====================================================================
% Appendix A --- Equivalences and overlaps of the closure matrix
% ---------------------------------------------------------------------
% New front appendix (2026-10). A short companion to Section 3: it states
% the overlaps the section discusses, gives the one-paragraph mechanism,
% and forward-points to the full proofs (moved to Appendix B, label
% app:proofs). It reuses the lemmas/propositions of that appendix BY \ref
% ONLY -- no label is duplicated here, so the proofs remain the single
% source. Kept to prose (~1.5 pages, double-spaced); the classification
% table already lives in Appendix B (Cor.~\ref{cor:class}).
% =====================================================================
\section{Equivalences and overlaps of the closure matrix}\label{app:equivalences}

\noindent Section~\ref{sec:two-questions} observes that several of the five
labour-market closures overlap at characteristic parameter values. This
appendix collects those overlaps, states why each holds, and marks what kind
of statement it is --- an exact identity, a limit, or a definitional
inclusion. The derivations are in Appendix~\ref{app:proofs}; only the
mechanism is kept here.

\subsection{Why overlaps arise}\label{app:eqwhy}

The programme enters final demand, so the shock is demand-only ($A \equiv
\one$) and touches neither the price nor the technology margin. Two features
of the kernel then make overlap the rule. First, the zero-profit conditions
are homogeneous of degree one in $(p,w)$ and contain no demand variable, so
relative prices are pinned by technology alone (Lemma~\ref{lem:price}) and,
under the consumption-price numeraire, the real wage sits at its baseline
anchor $w/\Pi(p) = 1$ for every demand vector, financing closure and shock
size (Corollary~\ref{cor:realwage}). A closure whose only task is to move the
\emph{wage} margin is therefore inert: the margin it is meant to clear is
already at rest. Second, at those prices the clearing system is exactly
affine and independent of $(\theta, \epsilon, \sigma)$
(Lemma~\ref{lem:affine}) --- the production side of the Samuelson
nonsubstitution theorem (Remark~\ref{rem:samuelson}) --- so any closure that
pegs the real wage solves the same affine system.

\subsection{The three pair-wise overlaps}\label{app:eqpairs}

\begin{itemize}
\item \textbf{Standard CGE $\equiv$ single-wage elastic supply.} The anchor
  binds, so the supply equation $\sum_i L_i = \Lbar\,(w/\Pi)^{\eta_s}$
  reduces to $\sum_i L_i = \Lbar$ for \emph{every} $\eta_s$: the two systems
  coincide equation by equation (Proposition~\ref{prop:ab}), and the
  single-wage elasticity $\eta_s$ is unidentified on the demand-only slice.
\item \textbf{Fixed real wage $\equiv$ Leontief closure.} With the real wage
  pegged, $(p, y) = (\one, y^{\ast})$ solves the system at every elasticity
  triple, and the Leontief corner selects the same candidate
  (Proposition~\ref{prop:gd}); the equivalence is exact, not an
  $\epsilon \to 0$ artefact.
\item \textbf{No-reallocation $\equiv$ the rigid corner of the sectoral
  family.} A frozen allocation and an infinitely inelastic supply curve are
  the same row of algebra: $\eta_{s,i} = 0$ for all $i$ reproduces the
  $\eta = 0$ endpoint (Proposition~\ref{prop:nest}).
\end{itemize}

\subsection{One family, two rigidities}\label{app:eqfamily}

The elastic closure is a curve rather than a point, and it connects the two
rigidities: at $\eta_s \to 0$ the allocation is frozen and the closure is the
no-reallocation case, while at $\eta_s \to \infty$ the supply slope diverges
and it approaches the fixed real wage --- a limit, never an attained cell.
The standard CGE closure is the zero-elasticity limit of the \emph{aggregate}
(single-wage) variant of the same family. The two limits differ in what the
vanishing elasticity freezes --- the sectoral allocation or only the total
--- which is why the family reaches both rigidities without coinciding with
either.

\subsection{Scope}\label{app:eqscope}

Every statement above is bounded by three conditions: the shock is
demand-only; the first equivalence uses the \emph{single-wage} form; and the
overlap concerns the price margin, since employment still adjusts under the
fixed-wage and Leontief closures. The sectoral family breaks the price-block
rigidity --- with $N$ wages, demand reaches prices through the wage structure
--- so the sectoral rows are distinct economies rather than further copies.
The breakdown off the demand-only slice, and the no-go result for nominal
wage rules, are derived in Section~\ref{sec:scope} of
Appendix~\ref{app:proofs}.